% --- Back cover ---
% Large PERSONIX wordmark, no logo image. Cream paper inherited globally.
% Absolute (tikz overlay) positioning so the footer sticks to the bottom.
\clearpage
\thispagestyle{empty}
\begin{tikzpicture}[remember picture, overlay]
  % --- Wordmark ---
  \node[anchor=north, font=\sffamily\bfseries\fontsize{48}{52}\selectfont]
    at ([yshift=-26mm]current page.north) {PERSONIX};
  \node[anchor=north, font=\rmfamily\itshape\large]
    at ([yshift=-44mm]current page.north) {Безкомпромісна зміна};
  \node[anchor=north, font=\rmfamily\small,
        text width=\paperwidth-50mm, align=center]
    at ([yshift=-54mm]current page.north)
    {Нецензуровна й непідкупна децентралізована репутаційна мережа};
  % --- Body ---
  \node[anchor=north, text width=\paperwidth-50mm, align=justify, font=\small]
    at ([yshift=-72mm]current page.north)
    {%
      Держава працювала, коли комунікація була повільною, рішення --- місцевими,
      а влада жила в тому самому місті, яким правила. Сьогоднішні мережі
      перевертають усі три припущення --- і збудовані на них інституції вже не
      пасують до світу, яким вони врядують. Ця книжка описує інструмент, який
      пасує: децентралізовану репутаційну мережу, у якій ідентичність
      залишається вашою, правда публічно аудитовна, а підкуп --- репутаційне
      самогубство.

      \smallskip
      Не маніфест. Не прогноз. Робоче креслення того, як можна збудувати
      наступника держави --- добровільно, одна декларація за раз.
    };
  % --- Footer (anchored to the bottom of the page) ---
  \node[anchor=south, font=\sffamily\footnotesize,
        text width=\paperwidth-50mm, align=center]
    at ([yshift=12mm]current page.south)
    {%
      Pavel Kudrna\quad\textbullet\quad Praha, 2026\par
      \href{https://personix.org}{personix.org}\par
      Ліцензовано на умовах Creative Commons Attribution-ShareAlike 4.0
      International (CC BY-SA 4.0).
    };
\end{tikzpicture}%
\mbox{}%
\clearpage
